\documentclass[10pt]{beamer}

% --- Theme Selection ---
\usetheme{default}
\usecolortheme{seagull}

% --- Packages ---
\usepackage[utf8]{inputenc}
\usepackage{graphicx}
\usepackage{booktabs} % For better looking tables

\usepackage[slovak]{babel}
\usepackage[T1]{fontenc}
\usepackage{lmodern} % or use \usepackage{libertine}

% --- Presentation Info ---
\title[1st Term Progress]{Lokálny terminálový asistent založený na LLM}
\subtitle{Bakalárska práca 1}
\author[Ostap Pelekh]{Autor: Ostap Pelekh\\Vedúci: Ing. Roderik Ploszek, PhD.}
\institute[Slovak University of Technology in Bratislava]{Faculty of electrical engineering and information technology (FEI)}
\date[]{}

\usefonttheme{professionalfonts} 

\begin{document}

% --- Slide 1: Title Page ---
\begin{frame}
    \titlepage
\end{frame}

% --- Slide 2: Research Question & Objectives ---
\begin{frame}{Cieľe práce}
    \begin{block}{Hlavný cieľ}
        Navrhnúť, implementovať a vyhodnotiť CLI/TUI asistenta, ktorý využíva lokálne LLM
        na interakciu s používateľom (z jednej strany) a so systémom (z druhej
        strany) s dôrazom na kontextovú citlivosť, a vyhľadávanie údajov z lokálneho
        i externého prostredia.
    \end{block}

\end{frame}

% --- Slide 3: Motivation ---
\begin{frame}{Motivácia: Prečo lokálne LLM?}
    \begin{itemize}
        \item \textbf{Bariéra syntaxe:} GenAI asistenti, využívajúc
        svoju schopnosť porozumieť prirodzenému jazyku a generovať odpovede i
        kódy, môžu zjednodušiť interakciu s terminálom a prebrať úlohy za
        používateľa.
        \item \textbf{Limity cloudových LLM:}
            \begin{itemize}
                \item \textbf{Súkromie:} Posielanie obsahu terminálu a histórie na servery tretích strán je bezpečnostné riziko.
                \item \textbf{Latencia a dosostupnosť:} Závislosť na internetovom pripojení spomaľuje workflow. API volania sú spoplatnené.
                \item \textbf{Kontrola nad modelom:} Obmedzené možnosti prispôsobenia a integrácie s lokálnymi nástrojmi.
            \end{itemize}
    \end{itemize}
    \vspace{0.3cm}
    \begin{center}
        \textbf{Riešenie:} Lokálny asistent bežiaci priamo na stroji používateľa.
    \end{center}

    % \vspace{0.3cm}
    % \scriptsize{ \textbf{Persónalná motivácia:} Takýto vyskumný projekt umožní
    % získať skúsenosti s prácou s LLM, pre ľahkú integráciou a optimalizáciou
    % pre špecifické úlohy neskôr.}
\end{frame}

% --- Slide 4: Research Question & Objectives ---
\begin{frame}{Cieľe práce}
    \textbf{Hlavné úlohy:}
    \begin{enumerate}
        \item Analýza existujúcich riešení, možností lokálnych modelov a možností integrácie s terminálom a externými nástrojmi.
        \item Návrh architektúry a modulov.
        \item Implementácia samostatných modulov.
        \item Vyhodnotenie úspešnosti a výkonu rôznych modelov pre rôzne prípady použitia a pri rôzných obmedzeniach zdrojov.
    \end{enumerate}
\end{frame}

% --- Slide 5. Progress to Date ---
\begin{frame}{Dosiahnuté výsledky k 1. semestru}
    \begin{itemize}
        \item[\checkmark] \textbf{Teoretická časť:} Spracovaná analýza teorie LLM, ich architektúry, organizácie a prepojenia s externými nástrojmi.
        \item[\checkmark] \textbf{Analytická časť:} Skušanie existujúcich CLI asistentov, skúmanie ich architektúr a riešení.
        \item[\checkmark] \textbf{Implementácia:} Naštudovanie technologií pre prácu s lokálnymi LLM.
        Komponent 1 (Klasifikátor) integrovaný do základného CLI.
    \end{itemize}
    \vspace{0.4cm}
    \begin{center}
        \textit{Aktualne prebieha:}
    \end{center}
    \begin{itemize}
        \item \textit{ ohodnotenie a výlešenie Komponentu 1 pre zvolené modely}
        \item \textit{ študium techník vyhľadávanie údajov a vývoj Komponentu 2.}
    \end{itemize}

\end{frame}

% --- Slide 5. Návrh a Komponent 1 ---
\begin{frame}{Architektúra a Komponent 1}
    % \textbf{Design Considerations:} Modularita medzi TUI a LLM backendom.
    \vspace{0.3cm}
    \begin{block}{Komponent 1: NL vs. Shell Command Classifier}
        \begin{itemize}
            \item Klasifikátor určený na rozlíšenie medzi \textbf{prirodzeným jazykom} (otázka) a \textbf{shell príkazom}.
            \item Skušanie prvého kroku k inteligentnému dopĺňaniu bez nutnosti manuálneho prepínania režimov.
            \item Adaptacia promptu pre zlepšenie presnosti klasifikácie.
        \end{itemize}
    \end{block}
    \begin{itemize}
        \item Integrácia cez \textbf{Ollama}.
        \item Skúšanie na \textbf{Qwen2.5(0.5b,1.5b), Gemma3(1b,4b)}
    \end{itemize}
\end{frame}


% --- Slide 6: Challenges \& Future Timeline ---
\begin{frame}{Úlohy a problémy do budúcnosti}
    \textbf{Výzvy:}
    \begin{itemize}
        \item Optimalizácia promptov pre rôzne modely a prípady použitia.
        \item Balansovanie výkonu modelu, využitia pamäte a latencie.
        \item Zvolenie architektúry pre efektívne retrieval z lokálnych a externých zdrojov.
        \item Technická krivka učenia pre MCP, LangChain, Finetuning techniques.
    \end{itemize}

    \vspace{0.3cm}
    \textbf{Plán na ďalšie mesiace:}
    \begin{table}[]
        \begin{tabular}{ll}
            \toprule
            \textbf{Month} & \textbf{Activity} \\ \midrule
            Februar        & Retrieval Mechanism + Komponent 2,3 \\
            Marec          & Testovanie Komponentov, Výlepšenia, Statistická Analýza \\
            April          & Dokumentácia, Finalizácia \\

            Po obhajobe    & Personalné vylepšenia, Finetuning, Výužitie nadobudnutých znalostí pri adaptácii modelov do rozličných prostredí. \\
        \end{tabular}
    \end{table}
\end{frame}

% --- Slide 7: Conclusion ---
\begin{frame}[fragile]{Záver}
    \begin{center}
        \begin{exampleblock}{\textit{"Let's make complicated things simple!"}}
            \texttt{\$ exit}
            \vspace{0.2cm}
        \end{exampleblock}
        \vspace{0.8cm}
        \textbf{Kontakt:} xpelekh@stuba.sk / xpelekh@is.stuba.sk
        % \textbf{GitHub:} ...?
    \end{center}
\end{frame}

\end{document}

Plan: 5min talk, 1m questions

% 30s Slide 1: Title Page ---
% 40s Slide 2: Main goal, tasks ---
% 40s Slide 3: Motivation ---
% 30s Slide 4. Progress to Date ---
% 30s Slide 5. Návrh a Komponent 1 ---
% 60s Slide 6: Challenges \& Future Timeline ---
% 60s Slide 7: Conclusion ---
